% Guidance and fine-tune paragraphs, gaps filled from the Re=1000 audit store and
% factual corrections applied. Key numbers: base bands [1,5)..[64,96) =
% 0.99/0.94/0.71/0.47/0.35; base+dial = 0.98/0.95/0.81/0.89/1.01; ft =
% 0.95/0.95/0.86/0.99/1.22; ft+dial = 0.94/0.95/0.92/1.06/1.01. MSE base 0.40,
% ft 0.44. Residual ratios: base 0.63, base+dial 0.42, ft 0.73, ft+dial 0.37.
% Corrected claim: the dial on the fine-tune does NOT improve low-k retention.

We first apply inference-time sampling guidance to the base model's denoising
trajectories. In Figure \ref{fig:spectrum-re1000} the energy spectrum curve of this
setup hides behind that of the ground truth for close to the entire range, while Figure
\ref{fig:spectrum-ratio-re1k} paints a slightly different picture. Table
\ref{tab:re-1k-ft-base} details that guidance raises the average energy content carried
by SR samples by $\approx10$, $40$ and $65$ percentage points over the base model across
the $k\in[16,32)$, $[32,64)$ and $[64,96)$ bands, recovering $80$ to $100\%$ of the
fine-scale energy carried over $k\in[16,96)$ by the ground truth samples. We note that
injecting fine-scale energy in this manner comes at two costs. First, it takes away
$1\%$ of the low-$k$ energy present in the $k\in[1,5)$ band, with a $9\%$ deficit at the
very largest scale $k=1$. Secondly, the guidance does not close the gap evenly across
the wavenumber range. The current formulation of the reward function pushes sampling
guidance to energise structures belonging to the $k\in[10,96)$ band evenly. Figure
\ref{fig:spectrum-ratio-re1k} shows the reward taking effect at $k=10$, where the
base$+$dial curve departs from the base curve, but it is not able to fully address the
mid-$k$ energy deficit, although it closes most of the high-$k$ energy gap over
$k\in[64,96)$.\newline

On the other hand, the finetuned model starts injecting energy as low as $k=7$ in its SR
samples, even though the reward formulation pushes over $k\in[10,96)$ specifically. It
registers retention levels as high as $86\%$ and $98.7\%$ over $k\in[16,32)$ and
$k\in[32,64)$, outperforming the base$+$dial generations under the same reward function.
It then over-energises structures with wavenumbers $k\in[41,113)$, peaking at
$1.24\times$ near $k=69$, before falling back under $\nicefrac{E(k)}{E_{GT}(k)}=1$ at
the end of the range. We note that it loses $5\%$ of the low-$k$ energy contained in the
large-scale structures of the flow, in exchange for the high-$k$ energy gains. Adding
the guidance on top of the finetuned model corrects the high-$k$ overshoot, bringing the
$[64,96)$ retention from $1.22$ back to $1.01$, and lifts the $[16,32)$ band further to
$0.92$, at the cost of pushing $[32,64)$ past the truth to $1.06$ and shaving another
percent off the low-$k$ energy, as seen in Table \ref{tab:re-1k-ft-base}. The plain
finetuned model achieves a $10\%$ worse MSE than the base model, but raises the mean
residual to $r_{x_0}=0.73\times r_{GT}$. The SR generations are slightly rougher, hence
more similar to the ground truth frames. We note that DDPO finetuning pushes samples
closer to the ground truth frames both in terms of energy and PDE residual, while
guidance only moves the energy content in the right direction despite being steered
towards the same objective, smoothing the residual further to $0.42\times r_{GT}$ on the
base model.\newline
